\documentclass[12pt,a4paper]{article}
\usepackage[utf8]{inputenc}
\usepackage{amsmath,amssymb,amsthm}
\usepackage{geometry}
\usepackage{algorithm}
\usepackage{algpseudocode}
\geometry{margin=1in}

\newtheorem{theorem}{Theorem}
\newtheorem{lemma}[theorem]{Lemma}
\newtheorem{proposition}[theorem]{Proposition}
\newtheorem{corollary}[theorem]{Corollary}
\theoremstyle{definition}
\newtheorem{definition}[theorem]{Definition}
\theoremstyle{remark}
\newtheorem*{remark}{Remark}

\title{Problem 57: Square Root Convergents}
\author{Project Euler Solution}
\date{}

\begin{document}
\maketitle

\section{Problem Statement}
The continued fraction for $\sqrt{2}$ is $[1; \overline{2}]$:
\[
\sqrt{2} = 1 + \cfrac{1}{2 + \cfrac{1}{2 + \cfrac{1}{2 + \cdots}}}
\]
In the first $1000$ expansions, how many fractions $p_n/q_n$ have a numerator with more digits than the denominator?

\section{Mathematical Development}

\begin{definition}[Expansion Sequence]
Define $p_0/q_0 = 1/1$ and the recurrence
\begin{align}
p_n &= p_{n-1} + 2q_{n-1}, \\
q_n &= p_{n-1} + q_{n-1},
\end{align}
for $n \geq 1$. The fraction $p_n/q_n$ is the $n$-th expansion of the continued fraction for $\sqrt{2}$.
\end{definition}

\begin{theorem}[Recurrence Derivation]
The $n$-th expansion satisfies $p_n/q_n = 1 + 1/(1 + p_{n-1}/q_{n-1})$, which yields the stated recurrence.
\end{theorem}

\begin{proof}
\[
1 + \frac{1}{1 + \frac{p_{n-1}}{q_{n-1}}} = 1 + \frac{q_{n-1}}{q_{n-1} + p_{n-1}} = \frac{p_{n-1} + 2q_{n-1}}{p_{n-1} + q_{n-1}}.
\]
Reading off numerator and denominator gives the recurrence. The base case $p_1/q_1 = 3/2$ matches $1 + 1/2$.
\end{proof}

\begin{theorem}[Matrix Formulation]
\[
\begin{pmatrix} p_n \\ q_n \end{pmatrix}
= A^n \begin{pmatrix} 1 \\ 1 \end{pmatrix}, \qquad
A = \begin{pmatrix} 1 & 2 \\ 1 & 1 \end{pmatrix}.
\]
\end{theorem}

\begin{proof}
Direct induction using $A \begin{pmatrix} p_{n-1} \\ q_{n-1} \end{pmatrix} = \begin{pmatrix} p_n \\ q_n \end{pmatrix}$.
\end{proof}

\begin{lemma}[Eigenvalues]
The matrix $A$ has eigenvalues $\lambda_1 = 1 + \sqrt{2}$ and $\lambda_2 = 1 - \sqrt{2}$.
\end{lemma}

\begin{proof}
The characteristic polynomial $\lambda^2 - 2\lambda - 1 = 0$ yields $\lambda = 1 \pm \sqrt{2}$.
\end{proof}

\begin{theorem}[Asymptotic Growth]
Both $p_n$ and $q_n$ grow as $\Theta\bigl((1+\sqrt{2})^n\bigr)$. The number of digits at step $n$ is
\[
D_n = \Theta\bigl(n \log_{10}(1+\sqrt{2})\bigr) \approx 0.3827\, n.
\]
Since $|\lambda_2| = \sqrt{2} - 1 < 1$, the $\lambda_2^n$ contribution decays exponentially, and $p_n/q_n \to \sqrt{2}$.
\end{theorem}

\begin{proof}
Decompose $\binom{1}{1}$ in the eigenbasis. The dominant eigenvalue $\lambda_1 = 1 + \sqrt{2}$ governs growth; $|\lambda_2|^n \to 0$ exponentially.
\end{proof}

\begin{lemma}[Pell Connection]
The denominators satisfy $q_n = 2q_{n-1} + q_{n-2}$, the recurrence for half-companion Pell numbers.
\end{lemma}

\begin{proof}
From $p_{n-1} = q_n - q_{n-1}$, substitute into the recurrence for $q_{n+1}$:
\[
q_{n+1} = p_n + q_n = (p_{n-1} + 2q_{n-1}) + q_n = (q_n - q_{n-1}) + 2q_{n-1} + q_n = 2q_n + q_{n-1}. \qedhere
\]
\end{proof}

\begin{theorem}[Digit Excess Criterion]
$p_n$ has more digits than $q_n$ if and only if $\lfloor \log_{10} p_n \rfloor > \lfloor \log_{10} q_n \rfloor$, equivalently, if there exists an integer $m$ with $q_n < 10^m \leq p_n$.
\end{theorem}

\begin{proof}
The digit count of $x$ is $\lfloor \log_{10} x \rfloor + 1$. A strict inequality holds iff a power of $10$ lies between $q_n$ (exclusive) and $p_n$ (inclusive).
\end{proof}

\section{Algorithm}

\begin{algorithm}[H]
\caption{Count Digit-Excess Convergents}
\begin{algorithmic}[1]
\State $p \gets 3$, $q \gets 2$, $\textit{count} \gets 0$
\For{$i = 1$ \textbf{to} $1000$}
    \If{$\text{digits}(p) > \text{digits}(q)$}
        \State $\textit{count} \gets \textit{count} + 1$
    \EndIf
    \State $(p, q) \gets (p + 2q,\; p + q)$
\EndFor
\State \Return \textit{count}
\end{algorithmic}
\end{algorithm}

\section{Complexity Analysis}

\begin{theorem}[Time Complexity]
The algorithm runs in $O(N^2)$ digit-level operations where $N = 1000$.
\end{theorem}

\begin{proof}
At step $n$, $p_n$ and $q_n$ have $O(n)$ digits. Each iteration performs two big-integer additions in $O(n)$ time. The total cost is $\sum_{n=1}^{N} O(n) = O(N^2)$.
\end{proof}

\noindent\textbf{Space:} $O(N)$ for storing two big integers with $\sim 383$ digits at step 1000.

\section{Result}
The answer is $\boxed{153}$.

\section{Answer}

\[
\boxed{153}
\]

\end{document}
